\begin{pastexam}[label={px:ncu-110-dsa-10}]{中央資工所 110 學年度 複選題第 10 題}
（複選；每題 5 分，全對才給分，答錯一個選項扣 1.5 分）

The adjacency list below is for an AOE network. The \textbf{end} field points to a list of
adjacent vertices, \textbf{dur} field is the duration of the activity, \textbf{link} field points to
another adjacent vertex, \textbf{vertex} field is the id of the adjacent vertex, \textbf{count} filed
is the number of immediate predecessors. Which of the following statements are true?
\begin{enumerate}[label=(\Alph*)]
\item The total duration of the critical path is 18.
\item Path 0, 1, 4, 6, 8 is not a critical path.
\item Path 0, 2, 4, 7, 8 is a critical path.
\item The latest time that event 4 (i.e., vertex 4) can occur is 8.
\end{enumerate}
\begin{center}
\begin{tikzpicture}[>={Stealth[length=2.2mm]},font=\small,scale=0.82,transform shape]
\node[font=\small\bfseries] at (0.30,0.62) {count};
\node[font=\small\bfseries] at (1.25,0.62) {end};
\node[font=\small\bfseries,anchor=west] at (2.60,0.62) {vertex dur link};
\draw (0,-0.31) rectangle (0.80,0.31); \draw (0.80,-0.31) rectangle (1.60,0.31);
\node at (-0.35,0.00) {0}; \node at (0.40,0.00) {0};
\draw[->] (1.20,0.00) -- (2.70,0.00);
\draw (2.70,-0.22) rectangle (4.80,0.22); \draw (3.40,-0.22) -- (3.40,0.22); \draw (4.10,-0.22) -- (4.10,0.22);
\node at (3.05,0.00) {1}; \node at (3.75,0.00) {5};
\draw[->] (4.45,0.00) -- (6.00,0.00);
\draw (6.00,-0.22) rectangle (8.10,0.22); \draw (6.70,-0.22) -- (6.70,0.22); \draw (7.40,-0.22) -- (7.40,0.22);
\node at (6.35,0.00) {2}; \node at (7.05,0.00) {4};
\draw[->] (7.75,0.00) -- (9.50,0.00);
\draw (9.50,-0.22) rectangle (11.60,0.22); \draw (10.20,-0.22) -- (10.20,0.22); \draw (10.90,-0.22) -- (10.90,0.22);
\node at (9.85,0.00) {3}; \node at (10.55,0.00) {5};
\node at (11.25,0.00) {0};
\draw (0,-0.93) rectangle (0.80,-0.31); \draw (0.80,-0.93) rectangle (1.60,-0.31);
\node at (-0.35,-0.62) {1}; \node at (0.40,-0.62) {1};
\draw[->] (1.20,-0.62) -- (2.70,-0.62);
\draw (2.70,-0.84) rectangle (4.80,-0.40); \draw (3.40,-0.84) -- (3.40,-0.40); \draw (4.10,-0.84) -- (4.10,-0.40);
\node at (3.05,-0.62) {4}; \node at (3.75,-0.62) {2};
\node at (4.45,-0.62) {0};
\draw (0,-1.55) rectangle (0.80,-0.93); \draw (0.80,-1.55) rectangle (1.60,-0.93);
\node at (-0.35,-1.24) {2}; \node at (0.40,-1.24) {1};
\draw[->] (1.20,-1.24) -- (2.70,-1.24);
\draw (2.70,-1.46) rectangle (4.80,-1.02); \draw (3.40,-1.46) -- (3.40,-1.02); \draw (4.10,-1.46) -- (4.10,-1.02);
\node at (3.05,-1.24) {4}; \node at (3.75,-1.24) {3};
\node at (4.45,-1.24) {0};
\draw (0,-2.17) rectangle (0.80,-1.55); \draw (0.80,-2.17) rectangle (1.60,-1.55);
\node at (-0.35,-1.86) {3}; \node at (0.40,-1.86) {1};
\draw[->] (1.20,-1.86) -- (2.70,-1.86);
\draw (2.70,-2.08) rectangle (4.80,-1.64); \draw (3.40,-2.08) -- (3.40,-1.64); \draw (4.10,-2.08) -- (4.10,-1.64);
\node at (3.05,-1.86) {5}; \node at (3.75,-1.86) {2};
\node at (4.45,-1.86) {0};
\draw (0,-2.79) rectangle (0.80,-2.17); \draw (0.80,-2.79) rectangle (1.60,-2.17);
\node at (-0.35,-2.48) {4}; \node at (0.40,-2.48) {2};
\draw[->] (1.20,-2.48) -- (2.70,-2.48);
\draw (2.70,-2.70) rectangle (4.80,-2.26); \draw (3.40,-2.70) -- (3.40,-2.26); \draw (4.10,-2.70) -- (4.10,-2.26);
\node at (3.05,-2.48) {6}; \node at (3.75,-2.48) {9};
\draw[->] (4.45,-2.48) -- (6.00,-2.48);
\draw (6.00,-2.70) rectangle (8.10,-2.26); \draw (6.70,-2.70) -- (6.70,-2.26); \draw (7.40,-2.70) -- (7.40,-2.26);
\node at (6.35,-2.48) {7}; \node at (7.05,-2.48) {7};
\node at (7.75,-2.48) {0};
\draw (0,-3.41) rectangle (0.80,-2.79); \draw (0.80,-3.41) rectangle (1.60,-2.79);
\node at (-0.35,-3.10) {5}; \node at (0.40,-3.10) {1};
\draw[->] (1.20,-3.10) -- (2.70,-3.10);
\draw (2.70,-3.32) rectangle (4.80,-2.88); \draw (3.40,-3.32) -- (3.40,-2.88); \draw (4.10,-3.32) -- (4.10,-2.88);
\node at (3.05,-3.10) {7}; \node at (3.75,-3.10) {5};
\node at (4.45,-3.10) {0};
\draw (0,-4.03) rectangle (0.80,-3.41); \draw (0.80,-4.03) rectangle (1.60,-3.41);
\node at (-0.35,-3.72) {6}; \node at (0.40,-3.72) {1};
\draw[->] (1.20,-3.72) -- (2.70,-3.72);
\draw (2.70,-3.94) rectangle (4.80,-3.50); \draw (3.40,-3.94) -- (3.40,-3.50); \draw (4.10,-3.94) -- (4.10,-3.50);
\node at (3.05,-3.72) {8}; \node at (3.75,-3.72) {2};
\node at (4.45,-3.72) {0};
\draw (0,-4.65) rectangle (0.80,-4.03); \draw (0.80,-4.65) rectangle (1.60,-4.03);
\node at (-0.35,-4.34) {7}; \node at (0.40,-4.34) {2};
\draw[->] (1.20,-4.34) -- (2.70,-4.34);
\draw (2.70,-4.56) rectangle (4.80,-4.12); \draw (3.40,-4.56) -- (3.40,-4.12); \draw (4.10,-4.56) -- (4.10,-4.12);
\node at (3.05,-4.34) {8}; \node at (3.75,-4.34) {5};
\node at (4.45,-4.34) {0};
\draw (0,-5.27) rectangle (0.80,-4.65); \draw (0.80,-5.27) rectangle (1.60,-4.65);
\node at (-0.35,-4.96) {8}; \node at (0.40,-4.96) {2};
\node at (1.20,-4.96) {0};
\end{tikzpicture}
\end{center}
\end{pastexam}

\begin{pastsoln}
答案：(B)(C)。

\textbf{讀圖}。 每個頂點的串列列出從它出發的活動（邊），格式是（終點，工期）；link 是 $0$ 的表示
串列結束。 共 $11$ 個活動：
\[
\begin{gathered}
0 \to 1\,(5),\quad 0 \to 2\,(4),\quad 0 \to 3\,(5),\quad 1 \to 4\,(2),\quad 2 \to 4\,(3),\quad 3 \to 5\,(2),\\
4 \to 6\,(9),\quad 4 \to 7\,(7),\quad 5 \to 7\,(5),\quad 6 \to 8\,(2),\quad 7 \to 8\,(5).
\end{gathered}
\]
count 欄和進入的邊數一致（$4$、$7$、$8$ 各有兩條進入的邊，$0$ 沒有），起點是 $0$、終點是 $8$。

\textbf{最早時間}（資結書第~22.5~節）。 依拓樸順序 $0, 1, 2, 3, 4, 5, 6, 7, 8$，
$\mathit{ee}[v] = \max_{u \to v}(\mathit{ee}[u] + \text{工期})$：
\[
\begin{aligned}
&\mathit{ee}[0] = 0,\ \mathit{ee}[1] = 5,\ \mathit{ee}[2] = 4,\ \mathit{ee}[3] = 5,\\
&\mathit{ee}[4] = \max(5 + 2, 4 + 3) = 7,\ \mathit{ee}[5] = 5 + 2 = 7,\ \mathit{ee}[6] = 7 + 9 = 16,\\
&\mathit{ee}[7] = \max(7 + 7, 7 + 5) = 14,\ \mathit{ee}[8] = \max(16 + 2, 14 + 5) = 19.
\end{aligned}
\]
整個工程最快 $19$ 完成，關鍵路徑長度是 $19$。

\textbf{最晚時間}。 從終點倒著算，$\mathit{le}[8] = 19$，$\mathit{le}[u] = \min_{u \to v}(\mathit{le}[v] - \text{工期})$：
\[
\begin{aligned}
&\mathit{le}[8] = 19,\ \mathit{le}[7] = 19 - 5 = 14,\ \mathit{le}[6] = 19 - 2 = 17,\ \mathit{le}[5] = 14 - 5 = 9,\\
&\mathit{le}[4] = \min(17 - 9, 14 - 7) = \min(8, 7) = 7,\ \mathit{le}[3] = 9 - 2 = 7,\\
&\mathit{le}[2] = 7 - 3 = 4,\ \mathit{le}[1] = 7 - 2 = 5,\ \mathit{le}[0] = \min(5 - 5, 4 - 4, 7 - 5) = 0.
\end{aligned}
\]

\textbf{關鍵活動}。 活動 $u \to v$ 的最早開始 $e = \mathit{ee}[u]$、最晚開始 $l = \mathit{le}[v] - \text{工期}$，
$e = l$ 的是關鍵活動：
\begin{center}
\small
\begin{tabular}{l*{11}{c}}
\toprule
活動 & $0{\to}1$ & $0{\to}2$ & $0{\to}3$ & $1{\to}4$ & $2{\to}4$ & $3{\to}5$ & $4{\to}6$ & $4{\to}7$ & $5{\to}7$ & $6{\to}8$ & $7{\to}8$ \\
\midrule
$e$ & 0 & 0 & 0 & 5 & 4 & 5 & 7 & 7 & 7 & 16 & 14 \\
$l$ & 0 & 0 & 2 & 5 & 4 & 7 & 8 & 7 & 9 & 17 & 14 \\
關鍵 & \checkmark & \checkmark & & \checkmark & \checkmark & & & \checkmark & & & \checkmark \\
\bottomrule
\end{tabular}
\end{center}
只由關鍵活動組成的從 $0$ 到 $8$ 的路徑有兩條：$0, 1, 4, 7, 8$（$5 + 2 + 7 + 5 = 19$）和
$0, 2, 4, 7, 8$（$4 + 3 + 7 + 5 = 19$），它們就是關鍵路徑。

\textbf{(A) 錯}。 關鍵路徑長度是 $19$，不是 $18$。 $18$ 是 $0, 1, 4, 6, 8$ 這類經過 $4 \to 6 \to 8$ 的路徑長。

\textbf{(B) 對}。 $0, 1, 4, 6, 8$ 的長度是 $5 + 2 + 9 + 2 = 18 < 19$，活動 $4 \to 6$、$6 \to 8$ 各有 $1$ 的
寬裕時間，不是關鍵路徑。

\textbf{(C) 對}。 $0, 2, 4, 7, 8$ 長度 $19$，每個活動都是關鍵活動，是關鍵路徑。

\textbf{(D) 錯}。 事件 $4$ 的最晚時間是 $\mathit{le}[4] = 7$，不是 $8$。 只看 $4 \to 6$ 那條邊會得到
$17 - 9 = 8$，但事件 $4$ 也是 $4 \to 7$ 的起點，要取最小值 $14 - 7 = 7$，否則 $4 \to 7 \to 8$ 會讓工程
延到 $20$。
\end{pastsoln}
